
@@M0@@, et si @@M1@@, @@M2@@. Ceci définit bien une filtration finie sur @@M3@@.

On définit un homomorphisme~:
@@M4@@
en posant @@M5@@ qui ne dépend visiblement pas du choix des @@M6@@ relevant les sections données @@M7@@ de @@M8@@. Si la suite exacte est localement scindée, @@M9@@ est un isomorphisme, car localement @@M10@@.

Dans le cas où @@M11@@ sont localement facteurs directs de Modules libres de type fini, les @@M12@@ sont localement facteurs directs de libres de type fini, donc les @@M13@@ également~; on peut donc écrire dans @@M14@@~:
@@M15@@, d'où l'additivité.

Soient @@M16@@, @@M17@@ deux topos commutativement annelés, et @@M18@@ un morphisme de topos annelés. On sait (Exp. IV, 2.7) que @@M19@@ définit un morphisme d'anneau sur les @@M20@@ correspondants, soit @@M21@@. Comme l'image inverse commute aux puissances extérieures, le morphisme @@M22@@ correspondant sur les @@M23@@ est un @@M24@@-homomorphisme pour la structure qui vient d'être définie.

Signalons deux cas particuliers :
\begin{itemize}[label=-]
\item Soit @@M25@@ un schéma~; @@M26@@ est ainsi muni d'une structure de pré-@@M27@@-anneau; c'est ce cas qui servira dans la suite du séminaire, et il peut être considéré comme « modèle » pour la théorie générale des 
